%% ── Abstract ─────────────────────────────────────────────
\chapter*{Abstract}
\addcontentsline{toc}{chapter}{Abstract}
\thispagestyle{plain}

Physics-based electrochemical simulation of lithium-ion battery packs requires resolving coupled, nonlinear partial differential equations across multiple spatial domains and temporal scales, ranging from sub-micron solid-diffusion within active electrode particles to pack-level thermal conduction across hundreds of cells. The computational cost of such models---particularly the Doyle--Fuller--Newman (DFN) class---has historically restricted rigorous electrochemical analysis to single-cell or small-pack configurations.

This thesis presents \VIBE{} (\textbf{Vi}rtual \textbf{B}attery \textbf{E}nvironment), a modular, GPU-accelerated Python framework for high-fidelity, coupled electrochemical--thermal--degradation simulation of large-scale heterogeneous battery packs. The framework makes four principal research contributions:

\begin{enumerate}[leftmargin=1.5em]
  \item \textbf{GPU-Accelerated Pack Scalability.} A generalised Kirchhoff-network current-distribution solver is tightly coupled to per-cell DFN-level physics. By leveraging PyTorch's batched tensor operations (\texttt{torch.vmap}) and reverse-mode automatic differentiation (\texttt{torch.func.jacrev}), the framework is shown to successfully scale large packs (e.g., 10S10P) with near-linear time complexity, breaking the exponential computational barrier of traditional CPU solvers.

  \item \textbf{Exact Local Sensitivity Analysis.} The fully differentiable tensor architecture enables exact local differential sensitivity analysis (e.g., SOC-dependent temperature rise sensitivity) directly through reverse-mode automatic differentiation, eliminating the truncation errors and massive $O(N)$ computational overhead of traditional finite-difference perturbation methods.

  \item \textbf{Coupled SEI Degradation Modelling.} Seven mechanistically distinct SEI growth models (including solvent-diffusion limited, reaction limited, EC-reaction limited, and tunnelling) are implemented in a plug-and-play registry. This enables comparative degradation tracking under varying physical conditions directly inside the coupled pack simulation.

  \item \textbf{Modular PDE Framework with Interchangeable Discretisations.} A symbolic expression abstract-syntax-tree (AST) is introduced, enabling conservation-law PDEs to be declared in natural mathematical notation and automatically discretised using Finite Difference (FDM), Finite Volume (FVM), or Chebyshev spectral collocation methods, selected independently per spatial domain.
\end{enumerate}

All governing equations are integrated via an implicit Backward-Euler Newton scheme. An adaptive timestepping scheme based on Richardson extrapolation local-truncation-error (LTE) estimation is implemented for the advanced solver mode.

Validation experiments demonstrate sub-millivolt agreement with PyBaMM reference solutions for single-cell discharge and identical precision for SEI growth models. Pack-level studies quantify the profound impact of thermal heterogeneity on current redistribution, temperature divergence, and localized capacity fade in a 10-series$\times$10-parallel pack configuration.

\vspace{0.5em}
\noindent\textbf{Keywords:} Lithium-ion battery; Doyle--Fuller--Newman model; SEI growth; Pack simulation; GPU acceleration; Thermal heterogeneity; Differentiable physics; Sensitivity analysis.

\newpage
